\manualchapter{Timing, polyphony and practical limits}{sec:operation}
The fixed master clock is 3,579,545\,Hz, matching the nominal clock documented
in ymfm's OPM phase-table derivation. The decoded YM3012 DAC runs at
clock/64 = 55,930.390625\,Hz. Key code/fraction uses 1/64 semitone steps
(1.5625 cents, at most 0.78125 cents rounding before chip table error).
Common pitch clamps from C\#0 through just below C\#8, approximately
$-3.9167$ to $4.0820$\,V relative to C4. Multipliers, DT1/DT2 and PM can
push operator frequencies higher; they retain chip aliasing and limits.

A 32-tap Blackman-windowed low-pass FIR at the native rate uses cutoff
min(20\,kHz, 0.4 times host rate), followed by fractional linear resampling.
This smooths DAC images; it is not an ideal brick-wall filter. The chip's own
high-frequency aliasing remains part of synthesis. Sample-rate changes preserve
chip and envelope state. The filter adds about 0.3\,ms of latency.

The control scan runs every 16 audio frames. Writes use one address/data
transaction per 64 master clocks, honoring the core's busy period. Replaceable
controls coalesce; pitch code/fraction pairs remain ordered and the shared
AM/PM depth register has independent pending values. Key edges are sampled
every audio frame and enter a 64-entry FIFO. Initial controls settle before
the first note. Regular key latency is at most two native sample periods
plus host sampling; a retrigger adds one further native sample period.
A burst can take up to roughly 1.2\,ms to drain, plus startup configuration.

Sustained retrigger rates up to 1\,kHz per lane are supported. Above the queue
limit a new rise/retrigger is rejected as a whole; one slot remains reserved
for release. Accepted key events are never coalesced. This is finite chip
bus bandwidth, not audio-rate arbitrary modulation. Continuous key traffic
outside the supported rate can delay control updates.

There are 1--16 lanes, determined by the widest input. Mono inputs broadcast;
missing channels of shorter polyphonic cables read zero. Disconnecting an
input removes its contribution. Removed lanes reset immediately, and return
with clean chip, LFO, noise and trigger state. Inactive lanes do not synthesize.
Reset and reload begin silently until gated; oscillator phase is not saved.

Decoded PCM is scaled by 10/32768 volts and output LEVEL, then limited to
$\pm10$\,V. A single full-level sine is approximately 2.5\,V peak before
LEVEL; four carriers can approach full-scale. Each routed jack carries its
own full-level signal, with no artificial stereo widening or output summing.
Unplugging R does not change L. Noise follows operator 4's envelope and VOL,
and is audible only when the chosen algorithm routes that operator to output.

\manualchapter{CPU and troubleshooting}{sec:practical}
The production engine is the pinned ymfm OPM core with mechanical C++11
compatibility changes. Nuked-OPM is retained only as an independent test
reference: the initial 16-chip prototype exceeded one thread's real-time
budget. On the local macOS arm64 validation machine, production Rack-module
benchmarks at 44.1, 48 and 96\,kHz with 64/256-frame workloads measured
roughly 10--12\% median callback time for 16 lanes, single-threaded.
These are headless production-module measurements, not a guarantee for every
machine or a complete patch. Leave headroom for other modules and audio I/O.

If silent, check GATE, output routing, carrier VOL and AR. If a held sound
fades, reduce SR or D1 and inspect SL. If AM does nothing, enable AM on an
audible operator and raise AMD and AMS. If a patch changes when a poly cable
shrinks, remember that shorter polyphonic cables supply zero on missing
lanes. Light/dark panels follow Rack's global preference without changing
sound or saved parameters.
